\documentclass[12pt,a4paper]{article}

% Paquetes básicos
\usepackage[spanish]{babel}
\usepackage[utf8]{inputenc}
\usepackage[T1]{fontenc}
\usepackage{lmodern}
\usepackage{setspace}             % Para interlineado
\usepackage{geometry}             % Para márgenes
\usepackage{fancyhdr}             % Encabezados y pies de página
\usepackage{graphicx}             % Inclusión de gráficos
\usepackage{array}                % Mejoras en tablas
\usepackage{booktabs}             % Tablas profesionales
\usepackage{longtable}            % Tablas que se extienden en varias páginas
\usepackage{hyperref}             % Hipervínculos
\usepackage{hyperref}
\hypersetup{
    colorlinks=true,    % Enlaces coloreados, sin recuadros
    linkcolor=black,     % Color de los enlaces internos
    urlcolor=blue,      % Color de las URL
    citecolor=blue,     % Color de las citas
    pdfborder={0 0 0}   % Quitar el borde del enlace
}
\usepackage{tikz}
\usepackage{eso-pic} % Permite agregar elementos a todas las páginas
\usetikzlibrary{calc} % Necesario para operaciones con coordenadas
% Configuración de márgenes
\geometry{
  top=2.5cm,
  bottom=2.5cm,
  left=3cm,
  right=2.5cm,
}
\usepackage{everypage}  % Permite ejecutar código en cada página
\usepackage{setspace}
\usepackage{booktabs}
\usepackage{caption}
\usepackage{capt-of} % Permite añadir captions fuera de un entorno float
\usepackage{float}
\usepackage{enumitem}
\usepackage{fontspec}
\setmainfont{Noto Sans} % O la fuente principal que prefieras
\newfontfamily\emoji{Noto Color Emoji} % Fuente que soporte emojis


\title{🔍 Dinámica: ``Detectives del Problema'' \\ \small{(Versión Extendida con Sistematización)}}
\author{Profesor: [Tu Nombre]}
\date{\today}

\begin{document}

\maketitle

\section*{🎯 Objetivo}
Aprender a identificar, delimitar y redactar adecuadamente un problema de investigación, e integrar un apartado de sistematización para explorar preguntas complementarias que enriquezcan el análisis.

\section*{👥 Modalidad y ⏱️ Duración}
\begin{itemize}[noitemsep]
    \item \textbf{Modalidad:} Virtual – Trabajo en grupos.
    \item \textbf{Duración:} 70 minutos (incluye tiempo adicional para la sistematización).
\end{itemize}

\section*{📌 Situación General}
En una universidad pública se ha observado una baja participación de los estudiantes en actividades de investigación durante los primeros años de carrera.

\section*{🧠 Dinámica de la Actividad}

\subsection*{I. Planteamiento del Problema (Ejemplo Resuelto – Grupo Modelo)}

\begin{tabular}{@{}p{5cm}p{10cm}@{}}
\toprule
\textbf{Elemento} & \textbf{Ejemplo del Grupo Modelo} \\ \midrule
\textbf{Tema General} & Participación estudiantil en investigación universitaria \\[1mm]
\textbf{Delimitación del Contexto} & Estudiantes de primer y segundo año de la Facultad de Ciencias de una universidad pública en Santo Domingo \\[1mm]
\textbf{Antecedentes} & Estudios breves indican que la falta de motivación y la percepción de que la investigación es exclusiva de estudiantes avanzados inciden en la baja participación. \\[1mm]
\textbf{Justificación} & Comprender las causas de la escasa vinculación temprana a la investigación es fundamental para diseñar estrategias que fortalezcan la formación científica desde los primeros años de carrera. \\[1mm]
\textbf{Relevancia y Pertinencia} & \begin{itemize}[noitemsep]
    \item \textbf{Relevancia:} Contribuye a mejorar la estructura curricular y la difusión de oportunidades investigativas.
    \item \textbf{Pertinencia:} Permite implementar cambios que optimicen la vinculación temprana en proyectos de investigación, generando un impacto positivo en la producción científica.
\end{itemize} \\[1mm]
\textbf{Causas del Problema} & \begin{itemize}[noitemsep]
    \item Falta de motivación.
    \item Poca difusión de oportunidades.
    \item Percepción de que la investigación es solo para estudiantes avanzados.
    \item Ausencia de asignaturas prácticas tempranas.
\end{itemize} \\[1mm]
\textbf{Consecuencias} & \begin{itemize}[noitemsep]
    \item Escasa formación investigativa en etapas tempranas.
    \item Menor producción científica estudiantil.
    \item Dificultad para cumplir con los objetivos curriculares.
\end{itemize} \\[1mm]
\textbf{Formulación del Problema} & ¿Cuáles son los factores que inciden en la baja participación en actividades de investigación de los estudiantes de primer y segundo año de la Facultad de Ciencias en una universidad pública en Santo Domingo? \\[1mm]
\textbf{Hipótesis} & La baja participación se debe a una combinación de falta de incentivos académicos, inadecuada difusión de oportunidades y una percepción negativa de la investigación. \\[1mm]
\textbf{Viabilidad} & Se cuenta con acceso a fuentes internas, la posibilidad de aplicar encuestas y entrevistas, y respaldo institucional para la recolección y análisis de datos. \\ \bottomrule
\end{tabular}

\subsection*{II. Apartado de Sistematización}

En este apartado se busca generar preguntas adicionales que desglosen y profundicen en distintos aspectos del problema. A continuación, se presenta un ejemplo:

\bigskip

\begin{tabular}{@{}p{5cm}p{10cm}@{}}
\toprule
\textbf{Aspecto a Investigar} & \textbf{Pregunta de Investigación Sistematizada} \\ \midrule
Factores Motivacionales & ¿Cómo influye la falta de incentivos académicos en la motivación de los estudiantes para participar en actividades de investigación? \\[2mm]
Difusión y Comunicación & ¿De qué manera afecta la comunicación institucional la percepción de las oportunidades de investigación entre los estudiantes? \\[2mm]
Percepción sobre la Investigación & ¿Cuál es la percepción de los estudiantes sobre la investigación en comparación con el aprendizaje teórico? \\[2mm]
Impacto Curricular & ¿Cómo se relaciona la estructura curricular de la Facultad de Ciencias con la baja participación en proyectos de investigación? \\ \bottomrule
\end{tabular}

\section*{💬 Estructura y Ejecución de la Dinámica en Clase}
\begin{enumerate}[label=\textbf{\arabic*.}, noitemsep]
    \item \textbf{Contextualización (10 min):} Explicar qué es un problema de investigación y la importancia de su sistematización, mostrando ejemplos de formulación y de preguntas sistematizadas.
    \item \textbf{Presentación de la Situación (5 min):} Plantear la situación general y consultar a los estudiantes si identifican otros problemas o aspectos relevantes.
    \item \textbf{Trabajo en Grupos: ``Detectives del Problema'' (25 min):} Cada grupo delimita el contexto, identifica causas y consecuencias, formula el problema y redacta una justificación breve utilizando la plantilla.
    \item \textbf{Incorporación de la Sistematización (20 min):} Una vez formulado el problema, cada grupo debe generar al menos tres preguntas de sistematización que profundicen en diferentes aspectos del problema.
    \item \textbf{Exposición y Retroalimentación (10 min):} Cada grupo comparte su formulación y sus preguntas sistematizadas. Se brinda retroalimentación colectiva, destacando buenas prácticas y sugerencias de mejora.
\end{enumerate}

\end{document}
